%%%%%%%%%%%%%%%%%%%%%%%%%%%%%%%%%%%%%%%%%%%%%%%%%%%%%%%%%
% SECCIÓN: RESUMEN Y ABSTRACT
%%%%%%%%%%%%%%%%%%%%%%%%%%%%%%%%%%%%%%%%%%%%%%%%%%%%%%%%%

\chapter*{}
\thispagestyle{empty}

\begin{center}
    {\large\bfseries Diseño e Implementación de un Simulador de Códigos BCH para el Análisis de Rendimiento en Canales con Ruido}\\[0.5cm]
    ZhiXiang Zhou Wang
\end{center}

\vspace{0.7cm}
\noindent{\textbf{Palabras clave}: Códigos BCH, Cuerpos de Galois, Simulación, Rendimiento, Teoría de la Información, Python.}

\vspace{0.7cm}
\noindent{\textbf{Resumen}}

Este Trabajo Fin de Grado presenta el diseño y la implementación de un simulador integral para códigos de bloque cíclicos BCH (Bose-Chaudhuri-Hocquenghem). El núcleo del proyecto consiste en el desarrollo de un motor aritmético basado en Cuerpos de Galois $GF(2^m)$ que permite realizar operaciones polinomiales eficientes para la codificación y decodificación de datos. Se implementa el algoritmo de Berlekamp-Massey para la localización de errores y la búsqueda de Chien para su corrección. El software desarrollado permite modelar canales de comunicación ruidosos y realizar un análisis comparativo de la tasa de error de bit (BER) frente a diferentes relaciones señal-ruido, proporcionando una herramienta visual y cuantitativa para evaluar la robustez de estos códigos en entornos de transmisión digital.

\cleardoublepage

\thispagestyle{empty}
\begin{center}
    {\large\bfseries Design and Implementation of a BCH Code Simulator for Performance Analysis in Noisy Channels}\\[0.5cm]
    ZhiXiang Zhou Wang
\end{center}

\vspace{0.7cm}
\noindent{\textbf{Keywords}: BCH Codes, Galois Fields, Simulation, Performance, Information Theory, Python.}

\vspace{0.7cm}
\noindent{\textbf{Abstract}}

This Thesis focuses on the design and implementation of a comprehensive simulator for BCH (Bose-Chaudhuri-Hocquenghem) cyclic block codes. The core of the project involves developing an arithmetic engine based on Galois Fields $GF(2^m)$ to perform efficient polynomial operations for data encoding and decoding. The Berlekamp-Massey algorithm is implemented for error localization along with the Chien search for error correction. The developed software enables the modeling of noisy communication channels and conducts a comparative analysis of the Bit Error Rate (BER) against various signal-to-noise ratios, providing a visual and quantitative tool to evaluate the robustness of these codes in digital transmission environments.

\cleardoublepage

%%%%%%%%%%%%%%%%%%%%%%%%%%%%%%%%%%%%%%%%%%%%%%%%%%%%%%%%%
% SECCIÓN: AUTORIZACIÓN DE BIBLIOTECA
%%%%%%%%%%%%%%%%%%%%%%%%%%%%%%%%%%%%%%%%%%%%%%%%%%%%%%%%%

\chapter*{}
\thispagestyle{empty}

\noindent\rule[-1ex]{\textwidth}{2pt}\\[4.5ex]

Yo, \textbf{ZhiXiang Zhou Wang}, alumno de la titulación \textbf{Grado en Ingeniería Informática} de la \textbf{Escuela Técnica Superior de Ingenierías Informática y de Telecomunicación de la Universidad de Granada}, con DNI \textbf{12807058Z}, autorizo la ubicación de la siguiente copia de mi Trabajo Fin de Grado en la biblioteca del centro para que pueda ser consultada por las personas que lo deseen. 

\vspace{5cm}

\noindent Fdo: ZhiXiang Zhou Wang

\vspace{2cm}

\begin{flushright}
    Granada, a \rule{1cm}{0.4pt} de Julio de 2026. [cite: 41]
\end{flushright}

\cleardoublepage

%%%%%%%%%%%%%%%%%%%%%%%%%%%%%%%%%%%%%%%%%%%%%%%%%%%%%%%%%
% SECCIÓN: CERTIFICACIÓN DE DIRECTORES
%%%%%%%%%%%%%%%%%%%%%%%%%%%%%%%%%%%%%%%%%%%%%%%%%%%%%%%%%

\chapter*{}
\thispagestyle{empty}

\noindent\rule[-1ex]{\textwidth}{2pt}\\[4.5ex]

D. \textbf{Jesús García Miranda}, Profesor del Área de Álgebra del Departamento de Álgebra de la Universidad de Granada. 

\vspace{0.5cm}

\textbf{Informa:} 

\vspace{0.5cm}

Que el presente trabajo, titulado \textit{\textbf{Diseño e Implementación de un Simulador de Códigos BCH para el Análisis de Rendimiento en Canales con Ruido}}, ha sido realizado bajo su supervisión por \textbf{ZhiXiang Zhou Wang}, y autorizo la defensa de dicho trabajo ante el tribunal que corresponda.

\vspace{1.5cm}

Y para que conste, expide y firma el presente informe en Granada, a \rule{1cm}{0.4pt} de Julio de 2026. 

\vspace{1cm}

\textbf{El director:} 

\vspace{4cm}

\noindent \textbf{Jesús García Miranda} 

\cleardoublepage

%%%%%%%%%%%%%%%%%%%%%%%%%%%%%%%%%%%%%%%%%%%%%%%%%%%%%%%%%
% SECCIÓN: AGRADECIMIENTOS
%%%%%%%%%%%%%%%%%%%%%%%%%%%%%%%%%%%%%%%%%%%%%%%%%%%%%%%%%

\chapter*{Agradecimientos}
\thispagestyle{empty}

\vspace{1cm}


\cleardoublepage